\chapter{总结与展望}

\section{研究工作总结}

本文围绕"面向路网层次结构建模的双曲表示学习方法"这一主题，针对现有欧氏空间路网表示方法在建模路网层次结构时的根本局限，提出了两个递进的双曲学习框架，从理论和实验两个层面验证了双曲几何在路网表示学习中的有效性。

第四章提出的双曲路网图神经网络（HRGNN）采用隐式层次建模范式，做出了以下贡献。

在方法层面，HRGNN 设计了首个专为交通路网道路类型分类任务定制的双曲 GNN，将庞加莱球模型中的双曲几何系统性地引入路网的节点嵌入学习。双通道注意力机制同时处理局部邻域和重要性引导 DFS 路径，在双曲空间中实现了对路网多尺度结构特性的协同建模。自适应融合策略根据路网特征动态调整两条路径的贡献，避免了固定权重策略的局限。

在实验层面，HRGNN 在城市内分类任务上相比欧氏基线提升 Micro-F1 8.4\%--9.1\%，在跨城市迁移学习中相比最优双曲基线提升 Macro-F1 13.0\%。覆盖 29 个国家 50 个城市的 $\delta$-双曲性分析揭示了城市规划模式与路网几何特性之间的关系，为数据集设计提供了量化依据。

第五章提出的基于蕴含锥的双曲层次路网表示模型（HHRoad）采用显式层次建模范式，做出了以下贡献。

在方法层面，HHRoad 首次将蕴含锥机制引入路网表示学习，对路网三层层次结构施加了显式的几何约束，将欧氏空间中软赋值矩阵的隐式层次表达升级为双曲空间中蕴含锥不等式的硬几何约束。组合蕴含学习范式成功地迁移至路网领域，建立了"路网功能区-局部区域-路段层次"之间的语义类比。蕴含损失、层次对比损失、结构重构损失和轨迹重构损失在 Lorentz 双曲空间中实现了多重监督信号的协同优化。

在实验层面，HHRoad 在 VecCity 基准的五城市、三任务上以综合平均排名第一（3.07）的成绩超越七种先进基线，在轨迹相似性检索任务上相比最强基线的 ACC@3 最高提升 9.4\%。

HRGNN 和 HHRoad 之间存在深度的联系，而非简单的并列。HRGNN 回答了"双曲空间是否有助于路网层次建模"这一基础问题，通过在道路类型分类任务上的显著提升验证了双曲几何的有效性。这一验证同时揭示了隐式范式的局限——层次包含关系未能在嵌入空间中得到显式表达，模型不构建多层次结构化表示。

HHRoad 在此基础上推进，回答了"如何在双曲空间中显式约束路网层次关系"的深入问题。蕴含锥机制将层次包含关系转化为可验证的几何条件，三层结构化表示使路网的功能-结构-路段层次得到了完整的几何化表达。从隐式到显式的递进，体现了对路网层次结构双曲建模认识的逐步深化。

两项工作的互补性还体现在方法论维度上。HRGNN 侧重路段级别的特征判别性学习，采用庞加莱球模型和有监督标签信号；HHRoad 侧重多层次结构的几何组织，采用 Lorentz 模型并结合无监督、弱监督和自监督信号。两者的结合覆盖了路网表示学习的核心问题——如何在合适的几何空间中，同时学习路段的局部语义特征和全局层次结构关系。

\section{路网层次结构双曲建模的理论启示}

本文的实验结果从多个角度验证了第三章的理论分析。HRGNN 的性能优势与数据集 $\delta$-双曲性的正相关关系（威尼斯 $\delta=3.0$ 时改善最大，雅典 $\delta=32.0$ 时改善较小但仍显著），直接证实了双曲空间对路网树状特性的适配。HHRoad 的嵌入可视化（功能区近原点、路段远离原点）定性验证了第三章关于"双曲距离天然编码层次深度"的理论论述。

更重要的是，两项工作在不同任务（分类 vs 表示学习）、不同双曲模型（庞加莱球 vs Lorentz）和不同城市集合上的一致有效性，表明双曲几何对路网层次建模的优势是稳健的，不依赖于特定的实验设置。

两项工作的对比揭示了两种范式的适用条件。隐式范式适合层次信息作为辅助特征的判别性任务，实现简洁且有效，但在需要精确层次组织的通用表示学习中表达力不足。显式范式适合层次结构建模为核心目标的任务，能够产出具有几何可解释性的多层次表示，但需要更多的先验信息（层次划分）和更复杂的训练目标。

这一发现对未来的方法设计具有指导意义：研究者可以根据目标任务的特性选择合适的范式，而非一刀切地采用某种方法。

本文提出的组合蕴含学习范式迁移至路网领域的成功，为跨领域方法迁移提供了一个有价值的范例。蕴含锥机制的适用性并不局限于路网——任何具有"通用-具体"层次结构的空间网络（如生物脉管网络、通信网络、供水管网）都可能从类似的双曲几何建模中受益。识别不同领域中结构类似的层次关系，并将双曲几何工具进行适配性迁移，是一种具有广泛潜力的方法论思路。

\section{研究局限性}

计算复杂度比较高。双曲运算（莫比乌斯加法、双曲距离计算、指数/对数映射）比欧氏对应操作计算代价更高，通常慢 2--3 倍。HHRoad 在北京（40,306 路段）上的每 epoch 训练时间约为 HRNR 的 2.3 倍。双曲空间的数值处理还需要额外的稳定性措施（梯度裁剪、数值截断），进一步增加了实现复杂度。

具有一定的数据依赖性。HHRoad 的功能区学习依赖轨迹数据，在轨迹数据稀疏的城市中可能降低学习效果。谱聚类的性能对路网规模敏感，在极大规模路网上可能需要近似方法。HRGNN 对节点特征质量有一定依赖，在 OSM 数据质量参差不齐的地区可能受到影响。

模型的可拓展性具有一定的局限性。当前模型使用固定或全局共享的曲率（HRGNN 使用可学习但全局共享的 $c$，HHRoad 固定 $\kappa=1.0$），而路网不同层次的"双曲程度"可能不同，多曲率模型可能进一步提升性能。层次结构的三层划分也是固定的，更灵活的层次划分（自适应深度、动态层次数）可能带来进一步提升。两项工作均假设路网为静态图，未考虑道路建设、封闭等导致拓扑动态变化的场景，也未建模路段属性的时间变化。

\section{未来研究方向}

路网拓扑因道路改造、事故封路等原因发生变化，路段属性（如实时速度、拥堵状态）也随时间变化。时态双曲图神经网络需要在保持层次几何组织的同时捕获路网的动态演化，这既是理论挑战（双曲空间上的时序建模），也是实践需求（实时路况感知）。将时间维度引入双曲模型，可以同时在空间层次维度和时间维度上建模路网状态，支持交通流预测\cite{jin2024stgnn,han2024bigst}、异常检测等时空预测任务。

多曲率模型与自适应几何。未来工作可以探索层次感知多曲率——路网不同层次使用不同曲率以适应各层次的几何需求，Bachmann 等\cite{bachmann2020kgcn} 在等曲率黎曼流形 GCN 框架下已为这一方向提供了部分理论依据。乘积空间模型将双曲空间与欧氏空间、球面空间的乘积作为嵌入空间，可以同时建模树状层次关系（双曲）、对称关系（球面）和平坦特征（欧氏）。Riemannian 自编码器可以学习路网的最优几何结构作为潜在变量。

更多下游任务的拓展。双曲路网表示可以扩展至路径规划（利用双曲嵌入的层次几何加速启发式搜索）、交通流预测（将双曲路段嵌入与时序模型结合）、异常检测（利用双曲空间中异常路段与正常路段的距离差异检测交通事件）以及城市功能区划（以功能区嵌入支持城市区域识别和用地分类）等任务。

双曲空间与其他几何空间的融合。球面几何适合建模环形或对称关系，路网中的环路结构具有球面几何特性。探索将双曲空间（层次关系）、球面空间（环形对称）和欧氏空间（平坦属性）融合的统一几何学习框架，可能为路网的多维度结构特性提供更全面的几何表达。图对比学习\cite{you2020graphcl,zhu2021gca} 中的自适应增强和多视图学习思路，可为此类多几何空间统一表示学习提供无监督或自监督训练目标。

可解释性的深入研究。双曲嵌入内置的几何可解释性为路网分析提供了独特工具。未来可以探索基于蕴含锥的层次因果分析（利用蕴含关系推断路网中的功能因果链）、基于双曲距离的路网相似性图谱（跨城市的路网结构相似性分析）以及面向城市规划决策者的可解释路网优化建议（将双曲几何洞察转化为实践指导）。

\section{研究展望}

路网表示学习在智能交通领域具有广阔的应用前景。本文的工作表明，双曲几何为路网层次结构建模提供了强大而有理论依据的几何框架。随着城市化的持续推进和智能交通系统的深入发展，对路网高质量表示的需求将持续增长。

在理论方面，双曲几何与图神经网络的融合仍有大量基础问题有待探索，包括双曲图神经网络的表达能力理论、双曲嵌入的泛化界和收敛保证、以及多几何空间学习的统一理论框架。在应用方面，高质量路网嵌入将推动智能交通系统在自动驾驶地图表示、实时导航优化、城市规划辅助决策等领域取得新的突破，并可能为跨模态城市理解（路网与遥感影像、社交媒体数据的融合）提供几何基础。在方法论方面，本文提出的组合蕴含学习范式的跨领域迁移成功，表明在具有层次结构的不同领域之间，双曲几何方法的思想和工具具有广泛的迁移潜力。

本文的研究揭示了双曲几何在路网表示学习中的独特价值，为后续研究奠定了方法和实验基础，期待未来的工作能够进一步拓展双曲几何在城市计算和智能交通领域的应用边界。